% === Begin Label Data ===
% ID: 560
% 难度星级: 
% 标签: 
% 备注: 
% 组卷引用次数: 0
% === End  Label Data ===

\begin{problem}{2020}{G}{全国II卷（文）}{18}{统计}
某沙漠地区经过治理，生态系统得到很大改善，野生动物数量有所增加。为调查该地区某种野生动物的数量，将其分成面积相近的 200 个地块，从这些地块中用简单随机抽样的方法选取 20 个作为样本，调查得到样本数据 $ (x_i,y_i) $（$ i=1,2,\cdots,20 $），其中 $ x_i $ 和 $ y_i $ 分别表示第 $ i $ 个样区的植物覆盖面积（单位：公顷）和这种野生动物的数量，并计算得 $ \displaystyle\sum_{i=1}^{20}x_i=60 $，$ \displaystyle\sum_{i=1}^{20}y_i=1200 $，$ \displaystyle\sum_{i=1}^{20}(x_i-\overline{x})^2=80 $，$ \displaystyle\sum_{i=1}^{20}(y_i-\overline{y})^2=9000 $，$ \displaystyle\sum_{i=1}^{20}(x_i-\overline{x})(y_i-\overline{y})=800 $.

（1）求该地区这种野生动物数量的估计值（这种野生动物数量的估计值等于样本中这种野生动物数量的平均数乘以地块数）；

（2）求样本 $ (x_i,y_i) $（$ i=1,2,\cdots,20 $）的相关系数（精确到 0.01）；

（3）根据现有资料和数据，各地块间植物覆盖面积差异较大，为提高样本的代表性以获得该地区这种野生动物数量更准确的估计，请给出一种你认为更合理的抽样方法，并说明理由。

附：相关系数 $ r=\dfrac{\displaystyle\sum_{i=1}^{n}(x_i-\overline{x})(y_i-\overline{y})}{\sqrt{\displaystyle\sum_{i=1}^{n}(x_i-\overline{x})^2}\sqrt{\displaystyle\sum_{i=1}^{n}(y_i-\overline{y})^2}} $，$ \sqrt{2}\approx 1.414 $.
\end{problem}

\begin{answer}

\end{answer}

\begin{solutions}

\end{solutions}